% Generated draft: requires mathematical and source review.
\documentclass[11pt]{article}
\usepackage[margin=25mm]{geometry}
\usepackage{fontspec}
\setmainfont{Latin Modern Roman}
\usepackage{amsmath,amssymb,mathtools,unicode-math}
\setmathfont{Latin Modern Math}
\usepackage{enumitem,needspace,xurl,hyperref}
\setlength{\parindent}{0pt}
\setlength{\parskip}{5pt}
\setlength{\emergencystretch}{3em}
\newcommand{\N}{\mathbb{N}}
\newcommand{\Z}{\mathbb{Z}}
\newcommand{\Q}{\mathbb{Q}}
\newcommand{\R}{\mathbb{R}}
\newcommand{\C}{\mathbb{C}}
\newcommand{\F}{\mathbb{F}}
\newcommand{\A}{\mathbb{A}}
\newcommand{\PP}{\mathbb{P}}
\newcommand{\E}{\mathbb{E}}
\newcommand{\eps}{\varepsilon}
\newcommand{\dd}{\,\mathrm{d}}
\begin{document}
\section*{UnsolvedMath ID: 30004251}
\textbf{Source problem code:} catalogue entry
\subsection{Definitions and mathematical statement}
The supplied record presents a research programme rather than a fully specified conjecture. Its questions depend on constructions and hypotheses that the record does not define.

Let \(X\) denote a toric geometric object and \(\mathcal F\) a sheaf in the intended setting. Write \(H^i(X,\mathcal F)\), for \(i\geq 0\), for its intended sheaf cohomology, and write \((g_j)_{j\in J}\) for the family called the distinguished generators. This notation is conditional: the record does not specify the category of \(X\), the allowed sheaves, the coefficient objects, the cohomology construction, or the construction and indexing of the distinguished generators.

The questions are:
\begin{enumerate}
\item Explain the mathematical meaning of the distinguished generators and determine the hypotheses under which they form a basis. If the intended cohomology is an ordinary module over a specified commutative ring \(R\), and \(g_j\in H^i(X,\mathcal F)\), the basis question asks when the map
\[
\bigoplus_{j\in J} R\longrightarrow H^i(X,\mathcal F),\qquad (a_j)_{j\in J}\longmapsto\sum_{j\in J}a_jg_j
\]
is an isomorphism. The sum is finite because its argument belongs to a direct sum. This module interpretation is conditional and does not specify a meaning of basis for blue coefficient objects.
\item Determine what additional structure these generators induce on the Cox ring of \(X\), whenever that ring is defined in the intended setting. The record does not specify the hypotheses for its construction, its grading conventions, or the proposed additional structure.
\item Determine whether, and how, the distinguished generators can be used for explicit sheaf cohomology computations. The admissible inputs, desired outputs, and scope of such computations remain unspecified.
\item Investigate the usefulness of blue sheaf cohomology for toroidal or logarithmic schemes, and formulate an extension to tropical schemes. The intended definitions of blue sheaf cohomology and tropical schemes, the applicable classes of objects, and the required properties of an extension are not supplied.
\end{enumerate}
The original statement asks about usefulness for toroidal or logarithmic schemes; the cleaned statement instead asks for an extension to all three settings. These are different demands, and the supplied record does not resolve their relationship. Consequently, no unconditional assertion with fixed domains and quantifiers can be extracted without further source clarification.
\subsection{Short English statement}
Explain the distinguished generators in toric sheaf cohomology, characterize when they form a basis, identify their effect on the Cox ring, and investigate explicit computations. Explore blue sheaf cohomology for toroidal and logarithmic schemes and its extension to tropical schemes. The relevant constructions and hypotheses require clarification; this is a statement-only research draft requiring human mathematical review.
\subsection{Sources}
\begin{itemize}
\item Local UnsolvedMath registry: supplied complete record with permanent ID \(30004251\), title \emph{Blue Sheaf Cohomology for Toric and Tropical Geometry}, and exact source problem code \texttt{catalogue entry}. The collection is \emph{Oberwolfach Reports Open Problems}. No individual proposer is supplied; the proposed year is recorded as 2019.
\item Supplied report citation: \emph{Toric Geometry (2020)}, Oberwolfach Reports, pp.\ 2671--2726, DOI \texttt{10.4171/owr/2019/43}. Source URL and literature-source link: \url{https://doi.org/10.4171/owr/2019/43}. The background describes Oberwolfach Reports as documenting mathematical workshops and their open-problem discussions.
\item Supplied literature-source reference: \emph{Cech cohomology over F1}, \url{https://doi.org/10.1016/j.jalgebra.2017.04.029}. No authors or more specific bibliographic details are supplied for this reference.
\item Supplied dated literature curation, checked on \textbf{2026-08-22}: the status is recorded as \texttt{partially\_solved}. The supplied assessment states that \v{C}ech blue sheaf cohomology is defined for important blueprint classes and refines toric cohomology, while the extracted questions concerning generators, applications, and tropical extension remain unresolved. This is supplied dated curation derived from the OWR record, not an independently verified literature assessment. Its time-sensitive status and cited claims require independent verification.
\item The registry labels the statement \texttt{corrected\_verified} and says that mathematical content was preserved while notation, typography, and sentence structure were repaired. This is a supplied editorial label, not independent mathematical verification of this draft.
\end{itemize}
\end{document}
